%%%%%%%%%%%%%%%%%%%%%%%%%%%%%%%%%%%%%%%%%%%%%%%%%%%%%%%%%%%%%%%%%%%%%%%%%%%
%%  appendixB.tex  –  Детали реализации и гиперпараметры
%%%%%%%%%%%%%%%%%%%%%%%%%%%%%%%%%%%%%%%%%%%%%%%%%%%%%%%%%%%%%%%%%%%%%%%%%%%

\chapter{Детали реализации и гиперпараметры}
\label{app:implementation}

В данном приложении документированы программная инфраструктура, конфигурации гиперпараметров для всех семи методов, процедуры предобработки данных, вычислительная инфраструктура и протоколы оценки, использованные на протяжении экспериментальной работы Главы~\ref{ch:experiments} (Экспериментальные результаты и валидация). Сводка гиперпараметров представлена в Таблице~3.3 Главы~\ref{ch:adaptation}; данное приложение предоставляет полные послойные спецификации.


%%=====================================================================
\section{Программная инфраструктура и зависимости}
\label{app:software}
%%=====================================================================

Весь экспериментальный конвейер реализован на Python~3.10 с PyTorch~2.1 в качестве основной среды глубокого обучения. Интегрирование ОДУ использует библиотеку \texttt{torchdiffeq} (версия~0.2.3), предоставляющую решатели Рунге--Кутты с адаптивным шагом и поддержкой метода сопряжённых чувствительностей. Компоненты графовых нейронных сетей построены на PyTorch Geometric (версия~2.4.0) для передачи сообщений, пакетирования графов и выборки окрестностей. Симуляция федеративного обучения использует Flower (версия~1.7) для коммуникации клиент-сервер и агрегации. Вычисления оптимального транспорта используют библиотеку POT (версия~0.9.1) с GPU-ускоренными итерациями Синкхорна. Учёт дифференциальной приватности использует Opacus (версия~1.4) для обрезки градиентов по отдельным образцам и отслеживания через моменты-бухгалтер. Генерация состязательных атак использует Foolbox (версия~3.3.3) и Adversarial Robustness Toolbox (версия~1.15). Слои байесовского вывода используют пользовательский модуль вариационного внимания, построенный на основе репараметризации Blundell и др. Интеграция LLM запрашивает OpenAI API (GPT-3.5-turbo) со структурированными промптами few-shot. Все случайные затравки фиксированы на 42 для воспроизводимости, эксперименты повторяются с пятью затравками для вычисления доверительных интервалов.


%%=====================================================================
\section{Конфигурации гиперпараметров}
\label{app:hyperparams}
%%=====================================================================

\subsection{Гиперпараметры CT-TGNN}
\label{app:hp_cttgnn}

\begin{table}[htbp]
\centering
\caption{Конфигурация гиперпараметров CT-TGNN.}
\label{tab:hp_cttgnn}
\begin{tabular}{lll}
\toprule
Параметр & Значение & Примечания \\
\midrule
Слои передачи сообщений    & 4     & Гетерогенные типоспецифичные \\
Скрытая размерность        & 256   & На слой \\
Размерность состояния      & 64    & Состояние ОДУ \\
Решатель ОДУ               & DOPRI5 & Адаптивный Рунге--Кутта \\
Относительный допуск       & $10^{-5}$ & Точность решателя \\
Абсолютный допуск          & $10^{-7}$ & Точность решателя \\
Пост. времени $(\tau_1,\tau_2,\tau_3)$ & $(1,60,3600)$\,с & Пакет, сессия, кампания \\
Размер окна TABN           & 32 шага & Окно скользящих статистик \\
Скорость обучения          & $2\times10^{-4}$ & AdamW \\
Затухание весов            & $10^{-5}$ & L2-регуляризация \\
Размер батча               & 128 графов & На GPU \\
Норма обрезки градиента    & 1.0   & Максимальная норма градиента \\
Эпохи обучения             & 200   & С ранней остановкой \\
Терпение ранней остановки  & 10    & Потери валидации \\
\bottomrule
\end{tabular}
\end{table}


\subsection{Гиперпараметры TripleE-TGNN}
\label{app:hp_triplee}

\begin{table}[htbp]
\centering
\caption{Конфигурация гиперпараметров TripleE-TGNN.}
\label{tab:hp_triplee}
\begin{tabular}{lll}
\toprule
Параметр & Значение & Примечания \\
\midrule
Слоёв GAT на кодировщик    & 3     & Сервис, трассировка, узел \\
Головы внимания            & 8     & Многоголовое внимание \\
Скрытая размерность        & 128   & На кодировщик \\
Окно темпорального вним.   & 16 шагов & Исторический контекст \\
Сила EWC $\lambda_{\mathrm{ewc}}$ & 100 & Катастрофическое забывание \\
Вес согласованности $\lambda_{\mathrm{con}}$ & 0.1 & KL-дивергенция \\
Вес потерь сервисов $\lambda_s$ & 0.4 & Многозадачный \\
Вес потерь трассировок $\lambda_r$  & 0.3 & Многозадачный \\
Вес потерь узлов $\lambda_n$   & 0.3 & Многозадачный \\
Выборки Фишера             & 500   & Для диагонали EWC \\
Скорость обучения          & $2\times10^{-4}$ & AdamW \\
Размер батча               & 64 графа & На GPU \\
\bottomrule
\end{tabular}
\end{table}


\subsection{Гиперпараметры FedLLM-API}
\label{app:hp_fedllm}

\begin{table}[htbp]
\centering
\caption{Конфигурация гиперпараметров FedLLM-API.}
\label{tab:hp_fedllm}
\begin{tabular}{lll}
\toprule
Параметр & Значение & Примечания \\
\midrule
Федеративных клиентов $K$  & 10    & Симулированные SOC \\
Раундов коммуникации $T$   & 100   & Агрегация сервера \\
Локальных эпох на раунд    & 5     & Обучение клиента \\
Доля усечения              & 20\%  & Верхнее и нижнее \\
Порог обрезки ДП $C$       & 1.0   & Градиент по образцу \\
Масштаб шума ДП $\sigma$   & 0.8   & Гауссовский механизм \\
Бюджет приватности $(\epsilon,\delta)$ & $(0.85,10^{-5})$ & На раунд \\
Регуляризация Синкхорна $\lambda$ & 0.1 & Энтропийная \\
Итерации Синкхорна         & 50    & Сходимость \\
Чувствительность Лапласа ОТ $\Delta_T$ & 0.01 & Транспортный план \\
Вес ансамбля LLM $\lambda_{\mathrm{LLM}}$ & 0.3 & Настроено по валидации \\
Примеров few-shot          & 5     & На категорию атак \\
Скорость обучения          & $5\times10^{-4}$ & Локальный AdamW \\
\bottomrule
\end{tabular}
\end{table}


\subsection{Гиперпараметры прямо совместимого обнаружения}
\label{app:hp_pqidps}

\begin{table}[htbp]
\centering
\caption{Конфигурация гиперпараметров прямо совместимого обнаружения.}
\label{tab:hp_pqidps}
\begin{tabular}{lll}
\toprule
Параметр & Значение & Примечания \\
\midrule
Слоёв MLP кодировщика      & 3     & На семейство протоколов \\
Скрытые размерности кодир.  & 256, 128, 64 & Убывающие \\
Активация                  & ReLU  & Все скрытые слои \\
Головы кросс-внимания      & 8     & Многоголовое слияние \\
Размерность внимания       & 64    & На голову \\
PGD состязат. $\epsilon$   & 0.1   & В пространстве признаков \\
Итерации PGD               & 10    & При обучении \\
Размер шага PGD            & 0.02  & На итерацию \\
Бины гистограммы тайминга  & 20    & Признаки Kyber \\
Скорость обучения          & $1\times10^{-3}$ & AdamW \\
Размер батча               & 256   & На GPU \\
\bottomrule
\end{tabular}
\end{table}


\subsection{Гиперпараметры MambaShield}
\label{app:hp_mamba}

\begin{table}[htbp]
\centering
\caption{Конфигурация гиперпараметров MambaShield.}
\label{tab:hp_mamba}
\begin{tabular}{lll}
\toprule
Параметр & Значение & Примечания \\
\midrule
Размерность состояния      & 64    & Скрытое состояние SSM \\
Размер свёрточного ядра    & 4     & Локальный контекст \\
Коэффициент расширения     & 2     & Внутренняя размерность \\
Блоков Mamba               & 6     & Последовательных \\
Размер ансамбля учителей $M$ & 5   & Непересекающиеся подмножества \\
Температура дистилляции    & 3.0   & Температура softmax \\
Нач. доля отравления $\rho_0$ & 0.05 & Начало учебного плана \\
Скорость роста отравл. $\alpha$ & $5\times10^{-4}$ & На итерацию \\
Макс. доля отравления $\rho_{\max}$ & 0.40 & Предел учебного плана \\
Нач. вес дистилляции $\lambda_0$ & 0.1 & Начало нарастания \\
Скорость нарастания дист. $\beta$ & $1\times10^{-3}$ & На итерацию \\
Априорное PAC-Байеса $\sigma_P$ & 1.0 & Гауссовское априорное \\
Доверие PAC-Байеса $\delta$ & 0.05 & Доверие границы \\
Скорость обучения          & $2\times10^{-4}$ & AdamW \\
\bottomrule
\end{tabular}
\end{table}


\subsection{Гиперпараметры стохастического трансформера}
\label{app:hp_stochastic}

\begin{table}[htbp]
\centering
\caption{Конфигурация гиперпараметров стохастического трансформера.}
\label{tab:hp_stochastic}
\begin{tabular}{lll}
\toprule
Параметр & Значение & Примечания \\
\midrule
Слоёв кодировщика          & 4     & Блоки трансформера \\
Головы внимания            & 8     & На слой \\
Скрытая размерность        & 256   & Размерность модели \\
Размерность FFN            & 1024  & Прямое распространение \\
Дропаут                    & 0.1   & Стандартный и вариационный \\
Прямых проходов МК $M$     & 50    & При выводе \\
Дисперсия априорного $\sigma_P^2$ & 1.0 & Изотропный гауссовский \\
Вес потерь $\lambda_1$ (задача) & 1.0 & Кросс-энтропия \\
Вес потерь $\lambda_2$ (KL)  & 0.1 & Регуляризация апостериорного \\
Вес потерь $\lambda_3$ (состязат.) & 0.5 & Состязательный \\
Вес потерь $\lambda_4$ (калибр.) & 0.2 & Штраф ECE \\
Вес потерь $\lambda_5$ (рег.)  & 0.01 & L2 затухание весов \\
Скорость подъёма возмущ. $\alpha_\psi$ & $5\times10^{-3}$ & Внутренний цикл \\
Температурное масштабирование $T$ & 2.3 & Пост-калибровка \\
Бины калибровки            & 15    & Вычисление ECE \\
Скорость обучения          & $2\times10^{-4}$ & AdamW \\
\bottomrule
\end{tabular}
\end{table}


\subsection{Гиперпараметры теоретико-игровой системы}
\label{app:hp_game}

\begin{table}[htbp]
\centering
\caption{Конфигурация гиперпараметров теоретико-игровой системы.}
\label{tab:hp_game}
\begin{tabular}{lll}
\toprule
Параметр & Значение & Примечания \\
\midrule
Конфигураций обнаруж. $|\calC|$ & 20 & Пространство действий \\
Раундов взаимодействия $T$ & 10{,}000 & Онлайн-обучение \\
Скор. обуч. МВ $\eta$      & $\sqrt{\log 20/10000}\approx0.017$ & Оптимальная \\
Порог неопределённости $\tau$ & 0.3 & Направление эксперту \\
Штраф неопределённости $\kappa$ & 0.5 & Дисконт полезности \\
Тайм-аут перечисл. опор Нэша & 60\,с & На экземпляр игры \\
Нормализация выигрышей     & $[0,1]$ & На раунд \\
\bottomrule
\end{tabular}
\end{table}


%%=====================================================================
\section{Конвейер предобработки данных}
\label{app:preprocessing}
%%=====================================================================

\subsection{Инженерия признаков}
\label{app:features}

Сырые файлы захвата пакетов обрабатываются в записи уровня потоков с использованием CICFlowMeter для CIC-IoT-2023 и CSE-CICIDS2018 и Argus для UNSW-NB15. Признаки потоков включают длительность, счётчики байтов (прямые и обратные), счётчики пакетов, статистики межпакетных интервалов (среднее, стандартное отклонение, минимум, максимум), счётчики флагов (SYN, FIN, RST, PSH, ACK, URG), статистики размера окна и индикаторы типа протокола. Все числовые признаки проходят $z$-нормализацию: $\tilde{x}_j = (x_j-\mu_j)/\sigma_j$, где $\mu_j$ и $\sigma_j$ вычисляются на обучающем множестве и последовательно применяются к валидационному и тестовому множествам.

Категориальные признаки (тип протокола, сервис, флаги состояния) проходят кодирование one-hot с последующей редукцией размерности PCA, сохраняющей 95\% объяснённой дисперсии. Это снижает закодированную размерность с нескольких сотен бинарных индикаторов до 15--30 главных компонент без значительной потери информации.

Для построения графа темпоральный срез графа создаётся каждые $\Delta t=10$~секунд. Узлы идентифицируются по IP-адресу, рёбра создаются для каждого наблюдаемого потока внутри окна. Признаки узлов агрегируют все признаки потоков, инцидентных узлу, через пулинг средним и максимумом.


\subsection{Аугментация данных}
\label{app:augmentation}

Аугментация графовых данных применяет три преобразования при обучении. Случайное удаление рёбер удаляет каждое ребро независимо с вероятностью~0,1, поощряя робастность к неполным наблюдениям сети. Возмущение признаков узлов добавляет гауссовский шум со стандартным отклонением~0,05 к нормализованным векторам признаков. Темпоральный джиттер сдвигает временные метки событий на $\calN(0,0.1^2)$ секунд, делая модель устойчивой к неточности синхронизации часов. Все аугментации применяются стохастически для каждого мини-батча и отключаются при валидации и тестировании.


%%=====================================================================
\section{Вычислительная инфраструктура}
\label{app:infrastructure}
%%=====================================================================

\subsection{Инфраструктура обучения}
\label{app:training_infra}

Обучение проводится на кластере из четырёх GPU NVIDIA Tesla V100 с 32\,ГБ памяти HBM2 каждый, соединённых через NVLink для мультиGPU параллелизма по данным. Хост-система предоставляет два CPU Intel Xeon Gold 6248 (40 ядер всего), 384\,ГБ ОЗУ DDR4 и 4\,ТБ NVMe SSD хранилище. Обучение со смешанной точностью через PyTorch AMP API (FP16 для прямого и обратного прохода, FP32 для обновления параметров) снижает потребление памяти приблизительно на сорок процентов, позволяя использовать батчи в 1,7$\times$ крупнее, чем при обучении FP32. Распределённый параллелизм по данным с бэкендом NCCL масштабирует обучение на четыре GPU с 92\% эффективностью линейного масштабирования.


\subsection{Вывод и симуляция развёртывания}
\label{app:inference_opt}

Симуляция развёртывания использует одиночный GPU NVIDIA Tesla P100 с 16\,ГБ памяти, отражая реалистичные ограничения периферийного развёртывания. Бенчмарки вывода используют модели, скомпилированные TorchScript, с размером батча 32 для измерений задержки и 512 для измерений пропускной способности. Все замеры времени представляют медиану ста независимых запусков с 95\% доверительными интервалами, вычисленными бутстрап-ресемплингом. Энергопотребление измеряется через NVIDIA System Management Interface (\texttt{nvidia-smi}) с опросом потребляемой мощности с интервалом 100\,мс при устойчивом выводе.


%%=====================================================================
\section{Протоколы оценки}
\label{app:eval_protocols}
%%=====================================================================

\subsection{Стратегия кросс-валидации}
\label{app:cv}

Стратифицированное разделение сохраняет пропорции классов в обучающей (70\%), валидационной (15\%) и тестовой (15\%) выборках. Для темпоральных наборов данных (CIC-IoT-2023 и CSE-CICIDS2018) разделение соблюдает хронологический порядок: первые 70\% записей формируют обучающую выборку, следующие 15\% --- валидационную, последние 15\% --- тестовую, предотвращая темпоральную утечку. Для федеративных экспериментов каждый симулированный клиент получает непересекающийся раздел обучающей выборки, а гетерогенность распределения меток вводится через распределение Дирихле с параметром концентрации $\alpha=0.5$.


\subsection{Тестирование статистической значимости}
\label{app:significance}

Все сравнения методов представляют парные $t$-тесты по пяти случайным затравкам с коррекцией Бонферрони по числу сравнений. Результат объявляется статистически значимым на уровне $\alpha=0.05$ после коррекции. Размеры эффекта представлены как $d$ Коэна для количественной оценки практической значимости наряду со статистической значимостью. Доверительные интервалы для точности, F1 и AUC вычисляются через интервал оценки Вильсона; доверительные интервалы для задержки и пропускной способности вычисляются через перцентильные интервалы бутстрапа с 10\,000 повторных выборок.


%%=====================================================================
\section{Доступность кода и данных}
\label{app:availability}
%%=====================================================================

Полный исходный код, файлы конфигурации и обученные контрольные точки моделей для всех семи методов доступны в следующем репозитории: \texttt{https://github.com/rogerpanel/robustidps.ai} (архивирован под DOI: 10.5281/zenodo.19129512). Репозиторий включает автоматизированные скрипты загрузки и предобработки данных, конвейеры обучения и оценки для каждого метода, файлы конфигурации гиперпараметров в формате YAML и блокноты Jupyter для воспроизведения всех рисунков и таблиц Глав~\ref{ch:experiments} и~\ref{ch:platform}. Наборы данных CIC-IoT-2023 и CSE-CICIDS2018 общедоступны от Канадского института кибербезопасности. Набор данных UNSW-NB15 доступен от Университета Нового Южного Уэльса. Наборы данных коллекции ICS3D доступны от их соответствующих поддерживающих организаций по академическим лицензиям.
